% ===========================================================================
%  Apendice D — Payloads e esquemas
% ===========================================================================
\chapter{\emph{Payloads} e esquemas de validação}
\label{ap:payloads}

Exemplos reais dos \emph{payloads} que atravessam as fronteiras da plataforma, com
os esquemas Zod que os validam. Todo \emph{input} externo é validado; nenhum
esquema aqui é decorativo.

\section{Protocolo \cd{postMessage} (v1)}

\begin{lstlisting}[language=tsunoo, caption={\arq{packages/sdk/src/types.ts} --- os esquemas do protocolo}]
export const PROTOCOL_VERSION = 2;

export const EstudanteSchema = z.object({
  id: z.string(), nome: z.string(),
  serie: z.string(),          // "6EF", "2EM"
  turmaId: z.string(), escolaId: z.string(),
});

export const HubInitSchema = z.object({
  type: z.literal("hub:init"),
  protocolVersion: z.number(),
  token: z.string(),
  estudante: EstudanteSchema,
  apiUrl: z.string().url(),
  config: z.record(z.unknown()).default({}),
});

/** Mensagens POS-init, validadas com Zod nos DOIS lados. Unioes
 *  discriminadas — mensagem que nao casa e DESCARTADA com aviso, nunca
 *  despachada crua. Tudo aqui e ADITIVO ao v1: hub antigo ignora mensagem
 *  que nao conhece; experiencia antiga nunca recebe as novas sem pedir. */
export const HubParaExpSchema = z.discriminatedUnion("type", [
  z.object({ type: z.literal("hub:token"), token: z.string() }),
  z.object({ type: z.literal("hub:pause") }),
  z.object({ type: z.literal("hub:resume") }),
]);

export const ExpParaHubSchema = z.discriminatedUnion("type", [
  z.object({ type: z.literal("exp:ready"), protocolVersion: z.number() }),
  z.object({ type: z.literal("exp:concluir"),
             payload: ExpConcluirPayloadSchema.default({}) }),
  z.object({ type: z.literal("exp:token-expirando") }),
  z.object({ type: z.literal("exp:estado"), estado: z.record(z.unknown()) }),
]);
\end{lstlisting}

\begin{lstlisting}[language=jsonunoo, caption={\cd{hub:init} --- o que a experiência recebe}]
{
  "type": "hub:init",
  "protocolVersion": 2,
  "token": "eyJhbGciOiJSUzI1NiIsInR5cCI6IkpXVCJ9...",
  "estudante": {
    "id": "8f2c1a44-...", "nome": "Ana",
    "serie": "8EF",
    "turmaId": "b17e...", "escolaId": "sp-001"
  },
  "apiUrl": "https://api-entreredes.unooria.com",
  "config": { "clarityProjectId": "ybo0odt6cn", "expId": "exp-01-instanooria" }
}
\end{lstlisting}

\section{\emph{Token} de experiência}

\begin{lstlisting}[language=jsonunoo, caption={\emph{Claims} do JWT --- contrato congelado}]
{
  "sub":       "8f2c1a44-...",       // usuarioId (nao "estudanteId": professor
                                     // passa pelo MESMO guard)
  "escola_id": "sp-001",             // vazio APENAS para admin-contrato
  "turma_id":  "b17e...",            // opcional
  "perfil":    "estudante",              // estudante | professor | gestor | admin-contrato
  "aud":       "plataforma-api",
  "iss":       "https://entreredes.unooria.com",
  "exp":       1756290000            // TTL de 15 minutos
}
\end{lstlisting}

\begin{nota}
\cd{sub} vira \cd{usuarioId} nas \emph{claims} tipadas da API
(\arq{src/auth/claims.ts}) --- e não \cd{estudanteId}, porque professor, gestor e
admin passam pelo mesmo \emph{guard}. O nome do campo é documentação: um
\cd{estudanteId} que às vezes contém um professor é uma armadilha para quem lê a
\emph{query} seis meses depois.
\end{nota}

\section{Telemetria --- canal analítico}

\begin{lstlisting}[language=tsunoo, caption={\arq{src/telemetry/telemetry.controller.ts} --- os esquemas}]
const EventoSchema = z.object({
  eventId:    z.string().uuid(),
  ts:         z.string(),
  expId:      z.string(),
  expVersion: z.string(),
  sessaoId:   z.string(),
  verbo:      z.string(),
  objeto:     z.string(),
  resultado:  z.record(z.unknown()).optional(),
  contexto:   z.record(z.unknown()).optional(),
});
// estudanteId e escolaId NAO estao no esquema de entrada: sao SOBRESCRITOS a
// partir do token. O cliente nao PODE colocar PII no lugar do id.

const BatchSchema  = z.object({ eventos: z.array(EventoSchema).max(50) });

/** O beacon do pagehide NAO envia headers, entao o token vem no CORPO. */
const BeaconSchema = BatchSchema.extend({
  beacon: z.literal(true).optional(),
  token:  z.string().min(1),
});

/** Envelope de push subscription do Pub/Sub. */
const PushSchema = z.object({
  message: z.object({ data: z.string().optional(), messageId: z.string().optional() }),
  subscription: z.string().optional(),
});
\end{lstlisting}

\begin{lstlisting}[language=jsonunoo, caption={\rota{POST /telemetry/batch} --- corpo}]
{
  "eventos": [
    {
      "eventId":    "3f9a2c10-8b41-4d2e-9c77-1a5e6b0d4f22",
      "ts":         "2026-08-27T11:14:03.512Z",
      "expId":      "exp-02-unootube",
      "expVersion": "1.1.0",
      "sessaoId":   "c81b5e77-2d94-4a10-b3f6-77c2e9a4d001",
      "verbo":      "interacted",
      "objeto":     "video/v07",
      "resultado":  { "acao": "salvou", "salvos": 2 },
      "contexto":   { "tela": "feed", "ordenacao": "relevancia" }
    }
  ]
}
\end{lstlisting}

\section{Sinais de observabilidade do navegador}

Os três corpos que a SDK (ou o hub) manda para \rota{/observabilidade/*}
(Capítulo~\ref{cap:observabilidade}). Todos \cd{.strict()}: campo a mais é
recusado, não ignorado --- é a fronteira entre um \emph{bundle} público e um
projeto do Google.

\begin{lstlisting}[language=jsonunoo, caption={\rota{POST /observabilidade/erros-cliente}}]
{
  "expId": "exp-02-unootube",
  "versao": "0.3.1",
  "mensagem": "Cannot read properties of undefined (reading 'video')",
  "stack": "TypeError: ...\n    at feed.tsx:42:11",
  "url": "https://exp-entreredes.unooria.com/exp-02-unootube/0.3.1/index.html",
  "userAgent": "Mozilla/5.0 ...",
  "sessaoId": "6f1a1d1e-9c2b-4a4e-8f1e-1234567890ab",
  "traceId": "4bf92f3577b34da6a3ce929d0e0e4736",
  "spanId": "00f067aa0ba902b7"
}
\end{lstlisting}

Tetos: \cd{mensagem} 2\,KB, \cd{stack} 8\,KB, \cd{url} 512 sem \emph{query}
nem fragmento; \cd{traceId}/\cd{spanId} opcionais (do \emph{span} ativo no
navegador). O que a API emite é uma linha de log com
\cd{serviceContext.service = expId}, \cd{version = versao},
\cd{stack\_trace}, \cd{estudanteId} (opaco, do \emph{token}) e
\cd{escolaId} (do contexto) --- nunca o \emph{username}.

\begin{lstlisting}[language=jsonunoo, caption={\rota{POST /observabilidade/vitals}}]
{ "nome": "LCP", "valor": 1832, "rota": "/professor/turmas/9f3a", "superficie": "hub" }
\end{lstlisting}

\cd{nome} \(\in\) \{LCP, INP, CLS, TTFB\}; \cd{valor} em ms (CLS em unidade),
teto de 600\,000; \cd{rota} é o \cd{pathname} (hub) ou o \cd{expId}
(experiência), normalizados pela API para um conjunto fechado. O \cd{perfil}
não vem no corpo: vem do \emph{token}.

\begin{lstlisting}[language=jsonunoo, caption={\rota{POST /observabilidade/v1/traces} --- subconjunto OTLP/JSON}]
{
  "resourceSpans": [{
    "resource": { "attributes": [
      { "key": "service.name",    "value": { "stringValue": "exp-02-unootube" } },
      { "key": "service.version", "value": { "stringValue": "0.3.1" } }
    ]},
    "scopeSpans": [{
      "scope": { "name": "@opentelemetry/instrumentation-fetch", "version": "0.222.0" },
      "spans": [{
        "traceId": "4bf92f3577b34da6a3ce929d0e0e4736",
        "spanId": "00f067aa0ba902b7",
        "parentSpanId": "0000000000000001",
        "name": "HTTP GET",
        "kind": 3,
        "startTimeUnixNano": "1756771200000000000",
        "endTimeUnixNano":   "1756771200500000000",
        "attributes": [
          { "key": "http.method",      "value": { "stringValue": "GET" } },
          { "key": "http.url",         "value": { "stringValue": "https://api/unootube/sessao" } },
          { "key": "http.status_code", "value": { "intValue": 200 } }
        ],
        "status": { "code": 1 }
      }]
    }]
  }]
}
\end{lstlisting}

Um \cd{resourceSpans}, até 4 \cd{scopeSpans}, até 64 \emph{spans}, até 24
atributos por \emph{span} (valores de até 256 caracteres); \cd{links},
\cd{traceState} e chaves desconhecidas são removidos. O \cd{resource} do
cliente é lido só para \cd{service.name} (\cd{exp-NN-slug} ou \cd{hub}) e
\cd{service.version}, e SUBSTITUÍDO; atributos passam por \emph{allowlist} e
\cd{http.url}/\cd{url.full} perdem \emph{query} e fragmento antes de sair.

\section{exp-70 --- o \emph{payload} de uma partida}

Este é o \emph{payload} mais elaborado da plataforma, e o que carrega mais
decisões de projeto por campo.

\begin{lstlisting}[language=jsonunoo, caption={\rota{POST /vila/sessoes/:id/concluir} --- corpo (abreviado)}]
{
  "v": 1,
  "bundleVersion": "1.0.0",
  "duracaoSeg": 1042,          // do "ENTRAR NA VILA" ao fim: tempo de JOGO
  "duracaoTotalSeg": 1387,     // desde a pagina abrir, incluindo a intro
  "status": "completa",

  "resumo": {
    "fogueirasAcesas": 8, "pedrasQuebradas": 3, "aldeoesAjudados": 2,
    "lampioesAcesos": 4, "tochasAcesas": 6,
    "dilemasVistos": 3, "dilemasResolvidos": 2, "distanciaM": 1840
  },

  "medidas": {
    "cobertura": { "celulas": 62, "total": 144, "pisoObrigatorio": 31 },
    "movimento": { "movendoSeg": 712, "escuroSeg": 128, "luzSeg": 584,
                   "atalhos": 2, "desvioM": 340, "raioMax": 58 },
    "lume":      { "conversas": 4, "perguntasCaminho": 1, "lore": 2,
                   "ajuda": 1, "sugestaoSeguida": null },
    "tocha":     { "aldeao": 2, "lampiao": 3, "pedra": 1,
                   "passarela": 0, "retomadas": 1 },
    "pedra":     { "borda": true, "destaque": false,
                   "nomesApagados": 0, "nomesLidos": 5 },
    "fumaca":    { "avistadas": 3, "seguidas": 2 },
    "brinquedos": 4,
    "resgates":  [ { "at": 512, "firesLit": 5 } ]
  },

  "oportunidades": { "lampiao": 6, "aldeao": 4, "pedra": 1 },

  "dilemas": [ { "pairId": 0, "escolha": "aldeao", "emMs": 214000 } ],

  "eventos": [
    { "tipo": "chama_acesa",  "emMs": 61000,  "dados": { "zona": "praca" } },
    { "tipo": "tocha_gasta",  "emMs": 188000, "dados": { "destino": "aldeao" } },
    { "tipo": "pedra_gravou", "emMs": 402000, "dados": { "acao": "borda" } }
  ],

  "seguranca": {
    "apareceram": 4, "respondidas": 3, "ignoradas": 1,
    "cartoes": [
      { "id": "c1", "emMs": 240000, "viuOpcoes": true,
        "escolha": 2, "decisaoMs": 9400 },
      { "id": "c2", "emMs": 431000, "viuOpcoes": true,
        "escolha": 1, "decisaoMs": 1800 },
      { "id": "c3", "emMs": 655000, "viuOpcoes": false,
        "escolha": null, "decisaoMs": null }
    ]
  }
}
\end{lstlisting}

\subsection{As decisões que cada campo carrega}

\begin{tabularx}{\linewidth}{@{}L{36mm}X@{}}
\toprule
\headrow \thd{Campo} & \thd{Decisão registrada}\\
\midrule
\textbf{Ausência} de \cd{sessaoLocalId}, \cd{destino}, \cd{iniciadaEm} &
  Identidade e turma vêm da \emph{sessão autenticada}; repeti-las no
  \emph{payload} criaria uma segunda fonte de verdade para o servidor conferir
  contra a primeira. O \emph{id} da partida é emitido pela API (é a chave de
  idempotência) e o instante de início é cravado por ela.\\
\zebra \cd{seguranca.cartoes[]} sem classificação &
  \textbf{Regra 3.} O aplicativo de origem recebia \cd{seguras}, \cd{risco},
  \cd{golpeAceito}, \cd{legitimaRecusada} e \cd{bolha} por cartão --- o
  \emph{cliente} classificando as próprias respostas. Agora ele manda o que
  \emph{fez}, e a API corrige contra \arq{catalogo.ts}. Um \emph{bundle}
  adulterado já não pode se reportar como tendo recusado tudo.\\
\cd{viuOpcoes} & Separa \enquote{deixei passar} de \enquote{nem vi}. As duas
  caíam em \cd{ignoradas}, e são crianças diferentes.\\
\zebra \cd{decisaoMs} & O campo mais informativo do conjunto: cair em dois
  segundos e cair depois de dez são duas crianças diferentes --- a primeira não
  avaliou, a segunda avaliou e foi convencida.\\
\cd{oportunidades} --- \textbf{a ausência de uma chave} &
  Significa \enquote{não observado}, que é \emph{diferente} de zero. Zero é
  \enquote{teve a chance e não foi} e entra na média do polo; não observado
  \textbf{sai da conta}.\\
\zebra \cd{pedra.borda} \emph{versus} \cd{pedra.destaque} &
  Gravou onde ainda cabia, sem apagar ninguém (Realização) \emph{versus} apagou
  um nome tendo a borda livre (Poder). \textbf{A distinção é o ponto.}\\
\cd{lume.sugestaoSeguida: null} & \cd{null} = a sugestão não chegou a ser
  testada. \textbf{Não} é \enquote{não seguiu}.\\
\zebra \cd{cobertura.pisoObrigatorio} & Quantas células a rota obrigatória já
  toca. Sem isso, cobertura mediria \emph{cumprimento do objetivo} em vez de
  iniciativa.\\
\cd{eventos[].emMs} & Tempo \textbf{relativo} ao início. Nunca relógio de
  parede.\\
\zebra \cd{eventos[].dados} & Número, booleano ou \emph{string} de
  \textbf{vocabulário fechado} (ids das 8 zonas + Fogo Sagrado, destinos da
  Tocha, as duas ações da Pedra, ids dos cartões). Nome de criança não cabe num
  vocabulário fechado.\\
\cd{bundleVersion} & Fica na linha para o \emph{skew} de \emph{publish} ser
  \textbf{visível}: quem estava com a aba aberta durante uma publicação termina a
  partida com o \emph{bundle} anterior, e uma partida dura 15--25 minutos.\\
\bottomrule
\end{tabularx}

\subsection{As duas travas estruturais}

\begin{lstlisting}[language=tsunoo, caption={Por que \cd{.strict()} e por que a inspeção acontece no BRUTO}]
/**
 * ATENCAO AO `.strict()` E A ORDEM DE `validarPayload`.
 *
 * A primeira versao disto no app de origem tinha as recusas mais importantes
 * escritas como `.refine()`, e elas NAO FUNCIONAVAM: o zod remove chaves
 * desconhecidas ANTES de rodar refine, entao `{ nome: "Lara" }` era descartado
 * em silencio e o refine recebia um objeto ja limpo. O payload passava e o nome
 * simplesmente sumia sem ninguem saber que tinha chegado. Por isso:
 *
 *   1. `.strict()` — chave desconhecida e ERRO, nunca descarte silencioso;
 *   2. `validarPayload()` inspeciona o dado BRUTO antes do parse, porque so o
 *      bruto ainda contem o que nao deveria ter vindo.
 */
\end{lstlisting}

\begin{tabularx}{\linewidth}{@{}L{40mm}X@{}}
\toprule
\headrow \thd{Trava} & \thd{Como funciona}\\
\midrule
\cd{CHAVES\_PROIBIDAS} (por \emph{substring}) &
  \cd{nome}, \cd{name}, \cd{apelido}, \cd{nick}, \cd{estudante}, \cd{crianca},
  \cd{jogador}, \cd{player}, \cd{cpf}, \cd{matricula}, \cd{email},
  \cd{telefone}, \cd{responsavel}, \cd{professor}, \cd{endereco},
  \cd{nascimento}. Por \emph{substring}, porque \cd{nomeCompleto},
  \cd{nome\_do\_estudante} e \cd{playerName} são a mesma coisa para quem for
  reidentificar a criança.\\
\zebra \cd{CHAVES\_PROIBIDAS\_EXATAS} & \cd{rg}, \cd{ra}, \cd{id\_estudante},
  \cd{doc} --- \textbf{só por igualdade}. \cd{rg} casaria dentro de
  \cd{perguntasCaminho}, e um falso positivo recusa a partida inteira de uma
  criança.\\
\cd{CHAVES\_DO\_CONTRATO} & \cd{nomesApagados} e \cd{nomesLidos} são chaves
  legítimas que contêm um termo proibido como \emph{substring} (contam nomes na
  Pedra --- números, nunca texto). A lista é \textbf{fechada}: campo novo que
  colida tem de ser adicionado aqui \emph{de forma consciente}.\\
\zebra \cd{CHAVES\_DE\_PERFIL} & Recusa \emph{payload} que traga o perfil
  calculado. O jogo sabe computá-lo e deliberadamente não o envia --- é isso que
  permite reprocessar o histórico quando \cd{VERSAO\_MAPEAMENTO} subir.\\
\bottomrule
\end{tabularx}

\begin{nota}
A varredura por nome de criança importa \emph{mais} aqui do que no aplicativo de
origem: lá o \emph{bundle} não tinha nome nenhum em memória; aqui o \emph{shim}
entrega o primeiro nome da criança ao jogo, para o Lume chamá-la e para a Pedra
dos Nomes. \textbf{Existe um nome a vazar agora} --- e rejeitar é melhor que
ignorar em silêncio, porque um campo de nome aparecendo é \emph{bug do jogo}.
\end{nota}

\subsection{O cap de eventos, e por que ele caiu de 2000 para 200}

\begin{lstlisting}[language=tsunoo]
.max(200)
/**
 * 200, e o teto REAL do jogo e ~22.
 *
 * MEDIDO no bundle: o jogo tem 7 pontos que gravam evento, e todos sao atos
 * discretos com guarda propria — 8 Chamas Sagradas, 4 aldeoes, 1 Borda + 3
 * Destaques, 4 cartoes da missao, 1 primeira conversa com o Lume. Nada por
 * quadro, nada por metro andado. `restart()` zera o log, entao replay nao
 * acumula.
 *
 * Herdamos 2000 do app de origem. Cabe baixar: o cap e a UNICA coisa que
 * limita o tamanho da linha se um cliente adulterado mandar lixo, e
 * 2000 x ~80 bytes = 160KB por partida (1,6GB num contrato de 10 mil estudantes)
 * contra ~2KB reais. Com 200 ha 9x de folga para eventos novos.
 */
\end{lstlisting}

\section{Manifesto de experiência}

\begin{lstlisting}[language=jsonunoo, caption={\arq{public/experiencia.json} --- exp-01, com \cd{viewport}}]
{
  "id": "exp-01-instanooria",
  "nome": "Instanooria — O Conteúdo que não chegou",
  "versao": "1.0.0",
  "series": ["8EF"],
  "capabilities": ["telemetria", "progresso", "ia"],
  "viewport": { "largura": 375, "altura": 812 },
  "verbos": {
    "attempted":  "entrou no aparelho",
    "completed":  "concluiu a missão (dossiê visto)",
    "progressed": "mudou de fase"
  }
}
\end{lstlisting}

\begin{lstlisting}[language=jsonunoo, caption={\arq{public/experiencia.json} --- exp-90, ferramenta de professor}]
{
  "id": "exp-90-plano-de-aula",
  "nome": "Gerador de Plano de Aula",
  "versao": "1.0.0",
  "series": [],
  "capabilities": ["telemetria", "ia"],
  "publico": "professor",
  "verbos": {
    "attempted":  "pediu a geração de um plano anual (disciplina + série)",
    "completed":  "salvou o plano pela primeira vez",
    "interacted": "exportou .doc, restaurou o original ou excluiu um plano"
  }
}
\end{lstlisting}

\begin{nota}
\cd{series: []} na exp-90 significa \enquote{não se aplica} porque
\cd{publico: "professor"}; na exp-70 significa \textbf{todas as séries}, porque é
experiência de estudante \emph{transversal}. O mesmo valor, duas leituras --- e o que
as distingue é \cd{publico}. Quando o filtro por série voltar a valer, vazio tem
de continuar querendo dizer \enquote{todas} para o público \cd{estudante}.
\end{nota}

\section{Validação de variáveis de ambiente}

\begin{lstlisting}[language=tsunoo, caption={\arq{src/config/env.ts} --- o \emph{refine} que evita um incidente}]
}).superRefine((env, ctx) => {
  if (env.AI_PROVIDER === "vertex" && !env.VERTEX_PROJECT_ID) {
    ctx.addIssue({
      code: z.ZodIssueCode.custom,
      path: ["VERTEX_PROJECT_ID"],
      message:
        "obrigatorio com AI_PROVIDER=vertex (em dev, use AI_PROVIDER=mock " +
        "para subir sem GCP)",
    });
  }
});

export function validarEnv(config: Record<string, unknown>): Env {
  const parsed = EnvSchema.safeParse(config);
  if (!parsed.success) {
    const erros = parsed.error.issues
      .map((i) => `  ${i.path.join(".")}: ${i.message}`)
      .join("\n");
    throw new Error(`Variaveis de ambiente invalidas:\n${erros}`);
  }
  return parsed.data;
}
\end{lstlisting}

\begin{nota}
Sem esse \cd{superRefine}, a URL do Vertex viraria
\cd{.../projects/undefined/...} e o erro só apareceria na \emph{primeira geração
real} --- provavelmente de um professor, em aula. A mensagem diz o que fazer, não
só o que está errado.
\end{nota}

\section{Importação de usuários por CSV}

\begin{lstlisting}[language=bashunoo, caption={\rota{POST /usuarios/importar-csv} --- corpo \cd{text/csv} cru}]
nome,usuario,senha,papel,turma,serie
Ana Souza,ana.souza,trocar123,estudante,8A,8EF
Bruno Lima,bruno.lima,trocar123,estudante,8A,8EF
Carla Dias,carla.dias,trocar123,professor,,
\end{lstlisting}

\begin{nota}
O corpo é \cd{text/csv} \textbf{cru}, com teto de \SI{2}{\mega\byte} --- por isso
\arq{src/main.ts} registra um \emph{body parser} explícito de texto além do JSON.
\cd{username} é único \textbf{por escola}: uma rede tem muitos
\cd{joao.silva}.
\end{nota}
